\begin{tikzpicture}[
  font=\scriptsize,
  tingkat/.style={draw, rounded corners, align=center, text width=3.05cm, minimum height=1.25cm},
  catatan/.style={align=center, text width=3.05cm},
  panah/.style={-{Stealth[length=2mm]}, thick}]
  \node[tingkat, fill=blue!8] (insiden) at (0,0)
    {\textbf{Insiden siber}\\Tim tanggap organisasi\\Kepala operasi TI};
  \node[tingkat, fill=orange!12] (krisis) at (4.3,0)
    {\textbf{Krisis organisasi}\\Istilah internal Medikara\\Direksi memimpin};
  \node[tingkat, fill=red!22] (nasional) at (8.6,0)
    {\textbf{Krisis Siber}\\Situasi darurat nasional\\Presiden menetapkan};
  \draw[panah] (insiden) -- (krisis);
  \draw[panah] (krisis) -- (nasional);
  \node[catatan, below=2mm of insiden]
    {Batasi, catat kronologi, nilai dampak. Lapor bila paling sedikit berisiko tinggi};
  \node[catatan, below=2mm of krisis]
    {Pemicu: RTO terlewat, lebih dari satu rumah sakit, atau risiko cedera berat};
  \node[catatan, below=2mm of nasional]
    {Bukan keputusan Medikara. Sebagai PSE, Medikara menindaklanjuti peringatan dini};
  \node[anchor=north west, align=left, text width=11.4cm] at (-1.55,-2.75)
    {{\tiny\mdseries Pemicu krisis organisasi adalah pengandaian Medikara. Definisi dan penetapan Krisis
    Siber mengikuti Peraturan BSSN 2/2024 Pasal 1 angka 3 dan Pasal 13 ayat (2).}};
\end{tikzpicture}
